% ============================================================
\section{Problem Formulation}
\label{sec:problem}
% ============================================================

\textbf{IoT traffic representation.}
We represent network traffic as a multivariate time series
$\mathbf{X} = (\mathbf{x}_1, \ldots, \mathbf{x}_T) \in \R^{T \times F}$,
where $T = 128$ timesteps correspond to a sliding window of flow records
and $F = 12$ features capture packet counts, byte volumes, inter-arrival
times, and flow duration (see Appendix A for the full feature list).

\textbf{Hard-negative attack.}
Let $\mathbf{X}_b$ be a benign sequence.  An adversarial sequence
$\mathbf{X}_a$ is a \emph{hard-negative attack} at stealth level $s$ if:
\begin{align}
  \label{eq:stealth}
  \text{CosSim}(\phi(\mathbf{X}_a),\, \phi(\mathbf{X}_b)) \geq s,
  \qquad s \in [0, 1],
\end{align}
where $\phi(\cdot) \in \R^{T \cdot F}$ denotes the flattened feature vector.
Intuitively, $s = 0.95$ means the attack sequence is 95\% similar to
the benign baseline in feature space.

\textbf{Attack patterns.}
We formalise four IoT attack patterns that embed adversarial signals
with low perceptibility:
\begin{align}
  P_{\text{slow}}(\mathbf{X}, s) &:
    x_{t,8} \leftarrow x_{t,8}\,(1 + (1-s) \cdot 0.03 \cdot t / T),
    \label{eq:slow} \\
  P_{\text{lotl}}(\mathbf{X}, s) &:
    x_{t,7} \leftarrow x_{t,7}\,(1 + (1-s) \cdot 1.15)
    \;\text{ if } t \bmod 32 = 0,
    \label{eq:lotl} \\
  P_{\text{beacon}}(\mathbf{X}, s) &:
    x_{t,8} \leftarrow x_{t,8}\,(1 - (1-s) \cdot 0.02)
    \;\text{ if } t \bmod 16 = 0,
    \label{eq:beacon} \\
  P_{\text{proto}}(\mathbf{X}, s) &:
    x_{t,10} \leftarrow x_{t,10} + (1-s) \cdot 0.05 \cdot \epsilon_t,
    \quad \epsilon_t \sim \mathcal{N}(0,1).
    \label{eq:proto}
\end{align}
These correspond to slow data exfiltration (gradual byte increase,
Eq.~\ref{eq:slow}), living-off-the-land mimicry (periodic micro-bursts,
Eq.~\ref{eq:lotl}), C2 beaconing (periodic dips, Eq.~\ref{eq:beacon}),
and protocol timing anomalies (timing jitter, Eq.~\ref{eq:proto}).
Feature indices follow the \dataset{} column ordering.

\textbf{Protocol compliance.}
A generated sequence must satisfy IoT protocol constraints
$\mathcal{C}_\pi$ for protocol $\pi \in \{\text{Modbus, MQTT, CoAP}\}$:
\begin{equation}
  \mathbf{X}_a \in \mathcal{F}_\pi \coloneqq
  \{\mathbf{X} : \mathcal{C}_\pi(\mathbf{X}) = \text{valid}\},
\end{equation}
where $\mathcal{C}_\pi$ checks constraints such as packet size ranges
(Modbus: 7--260 bytes), valid function codes, and inter-packet timing.

\textbf{Detection task.}
Given a test sequence $\mathbf{X}$, the detector $f_\theta$ outputs a
binary label $\hat{y} \in \{0, 1\}$ (0 = benign, 1 = attack) and an
anomaly score $\hat{s} \in [0,1]$.  The goal is to minimise the
false positive rate (FPR $= \text{FP}/(\text{FP}+\text{TN})$) while
maximising recall (detection rate) and F1 score, specifically on
hard-negative attacks where $s \geq 0.85$.
